\section{A blAST használata}

\subsection{A blAST programok fordítása}

A blAST az előbb bevezetett C4 nyelven belül megvalósított \gls{eDSL}.
Ez a {\tt blast.c4a} csomag felhasználásával lehet elérni.
Ennek megfelelően, a C4 eszközökön keresztül lehetséges a használata.
A fordításhoz a {\tt c4c} programmal állítható elő a fordított objektumfájl (object file).
Ezt a {\tt c4ld} linkerrel lehetséges futtatható fájlokká összeilleszteni.
Itt meg kell adni a konkrét rendszerre telepített LLVM és Clang könyvtárakat -- ezek puszta méretük miatt nincsenek belecsomagolva a blAST könyvtárba.
Egy lehetséges parancssorozatot lehet látni \aref{lst:blastcc}.~listában, de lehetséges, hogy az utolsó két kapcsolót (amelyek egy beágyazott LLD linker számára kerülnek átadásra) finomítani kell.
Az {\tt llvm-config} parancs ebben segítséget nyújthat.

\begin{lstlisting}[language=sh,caption={Egy blAST program fordításának lépései},label=lst:blastcc]
> c4c my-query.c4
> c4ld -o my-query -+ -lblast -larray my-query.o -lLLVM-22 -lclang-cpp
\end{lstlisting}

A blAST-el fordított verziónak egyeznie kell a linker paraméterben megadott verzióval.
Ha a két LLVM verzió kompatibilis működhet eltérő esetben is, de ez nem garantált, és nem valószínű, mert LLVM-nek nem stabil a bináris interfésze.

Röviden megjegyzendő, hogy folyamatban van egy {\tt c4i} LLVM \gls{JIT} alapú bináris, amelyikkel külön fordítás nélkül futtathatóak C4 alapú programok.
Sajnos a dolgozat írásáig azonban nem készült el.

\subsection{A blAST programok összetétele}

Egy blAST program kettő előzetes, és kettő lényegi lépésből áll.
Utóbbi kettőt akárhányszor végre lehet hajtani a program teljes élettartama alatt, így akár komplex, projektek között végzett ellenőrzéseket is lehet végrehajtani.

A legelső lépés, a {\tt blast.c4a} csomag használta.
Az első rendes függvényhívás előtt azonban inicializálni is kell a blAST könyvárat: ezzel igazából a Clang számára szükséges rendszerfájlok helyét kell megadni.
Alapértelmezetten az {\tt /usr} előtag alatt keresi ezeket a fájlokat, ha esetleg máshova került ezek telepítése, akkor azt kell megadni.
Egyes helyeken verziótól függő helyet kaptak (\pl {\tt /usr/lib/llvm/22}), így ilyenkor a megfelelő verzióhoz tartozó fejléceket kell megadni.
Ez egy globális tulajdonság, így programonként csak egyszer lehetséges beállítani, a {\tt blast\_init} függvény segítségével.
Példakód látható \aref{lst:blast-init}.~listában.
Az összes alábbi kódrészletek az itt említett inicializálást feltételezik, de külön nem jelenítik meg.

\begin{lstlisting}[caption={A blAST könyvár inicializálása},label=lst:blast-init]
use <blast.c4a>

# alapértélmzéshez blast_init with_defaults
blast_init resource_dir "/usr/lib/llvm/22/lib/clang/22"
\end{lstlisting}

Amint elő lett készítve a blAST könyvtár, lehet a tényleges kereséseket megvalósítani.
Ez két lépésből áll: a forráskódok megadása, majd a konkrét lekérdezések megírása.

\subsubsection*{Konkrét fájlok ellenőrzése}

A forráskódokat kétféle módon lehet a blAST számára átadni.
Egyik lehetőség, hogy konkrétan a keresni kívánt fájlokat adjuk meg.
Ilyenkor a fájl fordításához használandó kapcsolókat is meg kell adni, olyan formátumban, amelyet maga a {\tt clang} vagy {\tt clang++} parancssor képes lenne feldolgozni.
Ennek a motivációját az alábbi szakasz részletesen tárgyalja.
Szintaktikája magából a fájlból, majd egy \cfn nil/0 végű listából áll, amely a kapcsolókat tartalmazza.
Konkrét példa látható \aref{lst:blast-single}.~listában.
Egyéb kiemelkedő figyelmet érdemlő szituáció a más rendszerekre való kereszt-fordítás ({\it cross-compiling}), amelyben a konkrét rendszer viselkedése változtathatja a kód értelmezését (\pl egyes fejlécek lehet, hogy nem léteznek a célrendszeren).

\begin{lstlisting}[caption={Egyszeri források megadása},label=lst:blast-single]
let src_ast blast_build_ast_single "source.c"
  "--target=sparc-sun-netbsd",
  "-ansi",
  "-Dmakro",
  nil
\end{lstlisting}

\subsubsection*{Fordítási adatbázis}\label{subsec:compdb}

A Clang könyvtárakban a fordítási adatbázis ({\it compilation database}, az alapértelmezett esetben egy {\tt compile\_commands.json} nevű JSON fájl formájában) látja el azt a feladatot, hogy a feldolgozó eszköz tudja, milyen paraméterekkel került fordításra egy-egy adott projekt fájlja, hogy helyesen tudja azt értelmezni.

A C és C++ nyelvekben ennek nagyon nagy jelentősége van, akár azt is el tudja dönteni, hogy egyáltalán melyik fájl számít szintaktigailag helyesnek.
Egy nagyon egyszerű példakód látható \aref{lst:modbrk}.~listában.
Bár két rendes sorból áll, legalább három féle módon lehetséges értelmezni a fordítási paraméterek függvényében.
Látszik alább, hogy egyik példa az előfeldolgozón keresztül módosítja a forrás értelmezését, ami alapvetően nem meglepő.
De a másik kettő példától az is látszik, hogy erre nincs feltétlen szükség.
A bemutatott példán kívül is lehet még sok ilyen helyzetet előállítani, például architektúra függő beépített függvényekkel.

\begin{lstlisting}[language=c++,label=lst:modbrk,caption={Egyszerű kódrészlet egy module struktúrával ({\tt modbrk.cxx})}]
struct module { };

module valami;
\end{lstlisting}

\begin{description}
\item[C++20 előtti verziókkal]
Minden {\tt -std=c++\bf 17} formájú kapcsoló, amelyik C++17 vagy annál öregebb C++ verziót határoz meg, sikeresen fordítja a kódot.
\begin{lstlisting}[language=sh]
> clang++ -c -std=c++17 modbrk.cxx
\end{lstlisting}

\item[C++20 és utána következő verziókkal]
A C++20-ban bevezetett {\tt module} kontextus függő kulcsszó miatt ez a kód már így nem fordul.
\begin{lstlisting}[language=sh]
> clang++ -c -std=c++20 modbrk.cxx
modbrk.cxx:3:1: error: module declaration must occur at the start of the translation unit
    3 | module valami;
      | ^
modbrk.cxx:1:1: note: add 'module;' to the start of the file to introduce a global module fragment
    1 | struct module { };
      | ^
modbrk.cxx:3:8: fatal error: module 'valami' not found
    3 | module valami;
      |        ^~~~~~~
2 errors generated.
\end{lstlisting}

\item[Makró felüldefiniálásával]
Bármelyik verzióban fordíthatóvá tehető, ha a fordító előtt az előfeldolgozó (preprocessor) szimplán átnevezi a problémás {\tt module} szót.
\begin{lstlisting}[language=sh]
> clang++ -c -std=c++20 -Dmodule=valami_mas modbrk.cxx
\end{lstlisting}
\end{description}

Ezen rugalmasság miatt van szüksége minden adott fájl feldolgozásához a teljes parancssorhoz, amivel a tényleges projekt fordítása során kerül fordításra.
Az előző szekcióban ezért kellett megadni a teljes kapcsolólistát, általában pedig a fordítási adatbázisra ezért van szükség.

Az ilyen adatbázisokat a blAST számára direkt módon át lehet adni, és akkor automatikusan az összes forrásfájlt ellenőrzi a program számunkra.
Ez lehetővé teszi a teljes projekteken keresztül végrehajtott kereséseket extra konfiguráció vagy kód nélkül.
Az ellenőrzött fájlok szűréssel előállított vagy bármi más módon vett részhalmazuk megadása tényleges ellenőrzésre jelenleg még nem támogatott.
Ennek a szintaktikája látható \aref{lst:blast-compdb}.~listán, bár nincs hozzá külön saját \gls{DSL} jellegű forma, csak egy egyszerű függvényhívás.

Érdemes itt kiemelni, hogy a C4 nyelv kiértékelése lusta.
Emiatt a fordítási adatbázis megadása azonnal nem értékeli ki azt, csak amikor -- ha egyáltalán -- ténylegesen használatba kerül az.
Ezt fel lehet használni arra, hogy sok különböző fordítási adatbázist támogatunk, és a kód később dönti csak el, hogy melyiket használja valójában.
Ilyen lehet, \pl egy ``Debug'' és egy ``Release'' verzió, amelyiket utána valami környezeti változó, vagy hasonló futásidejű érték alapján dönt el, hogy melyiken futtassa le, anélkül, hogy definiáláshoz sok különböző értéket kellene megadni.

\begin{lstlisting}[caption={Fordítási adatbázis megadása},label=lst:blast-compdb]
let ast blast_build_from_compile_commands "compile_commands.json"

let rel_ast blast_build_from_compile_commands "release/compile_commands.json"
let dbg_ast blast_build_from_compile_commands "debug/compile_commands.json"

# "when c t else f" a C-ből ismert "c ? t : f" operátorral azonos viselkedésű
let used_ast when (env "CI" == 1) rel_ast else dbg_ast
\end{lstlisting}

\subsection{Lekérdezések blAST programokban}

A lekérdezések képzik a legfontosabb részét a blAST alapú programoknak, ezek ebben a szakaszban kerülnek ismertetésre.
Általánosan egy SQL nyelvhez hasonló szintaktikájuk van, amelyik lehetővé teszi, hogy maguk a lekérdezések deklaratív módon a lényegi részekre fokuszáljanak.

Minden lekérdezések a \cfn blast/2 függvénnyel kezdődik és az \cfn end\_blast/0 függvényig tart.
A kettő között a blAST lekérdezések \gls{eDSL} szintaktikáját várják.
Ennek első paramétere, azon \gls{AST} érték, amelyeket az előző szekciók alapján a program definiált.
Ezt követik az keresések és az extra találatokat lekezelő függvények.
Az alábbi bekezdések ezeket írják le, majd \aref{lst:complex}.~listában látható egy összesített, komplexebb példa, amelyiken láthatóak ez egyes szintaktikai darabok összeszedve.

\subsubsection*{Keresések}

A blAST programok keresései elsődlegesen a {,,{\tt match into}''} függvényekkel kezdődnek, amelyet egy találatot kezelő függvény törzse követ.
Ezekről alább lesz részletesen szó.
A kezelő függvényt követi a tényleges lekérdezés kifejezése: itt kerül megadásra, milyen elemeket kell visszaadni az előbb definiált kezelő függvénynek.
Alapvetően három fő lehetőség van: vagy egy kifejezést keres az \cfn expression/0 függvénnyel, egy parancsot a \cfn statement/0-el, vagy egy deklarációt a \cfn declaration/0 megadásával.
Megjegyzendő, hogy a definíció is egy deklaráció fajta.

\begin{lstlisting}[caption={Különböző lekérdezések blAST-al},label={lst:queries}]
blast used_ast
  match into # <...>
  any declaration

  match into # <...>
  such declaration
  that is named "foo"
end_blast
\end{lstlisting}

Azt, hogy konkrétan milyen feltételek mentén kell ezeket megkeresni -- vagy akár az összes adott csúcsot az \gls{AST}-ben -- azt két szintaxist képző függvény adja meg: a \cfn any/1 és a ,,{\tt such/2 \rm ... \tt that/1 \rm ...}'' pár.
Ezek nevük alapján elvárt módon viselkednek: az \cfn any/1 után álló típusból az összes átadásra kerül a kezelő függvénynek, míg a \cfn such/2 {\it fajta} \cfn that/1 {\it feltétel}' konstrukció esetén, csak a megadott {\it feltétel} által ki nem szűrt elemek kerülnek tovább.

Szűrni azonban nem csak a \cfn such/2 függvénnyel lehetséges: a konkrétan keresett elem alfajtáját a megfelelő fő fajtát választó függvény hívásakor kell megadni, nem utólagosan.
Tehát, ha nem bármilyen deklarációt, csak függvény deklarációt kell megkeresni, azt a \cfn declaration/0 függvénynél kell megadni.
Példa az ilyen módon megadott specializált lekérdezésekre látható alább.

\begin{lstlisting}
blast used_ast
  match into # <...>
  any function declaration

  match into # <...>
  any call expression
end_blast
\end{lstlisting}

A szűrésben a \cfn that/1 függvény után egy olyan kifejezést lehet megadni, amelyik a az előzőek szerint megadott típusú \gls{AST} elemei közül választja ki, a keresetteket.
Ehhez számos függvény érhető el, \pl \cfn static/0, \cfn expanded/1, \cfn named/1.
A teljes listát nem sorolom fel, de az ASTMatcher konstrukcióknak megfelelő függvények elérhetőek.

Van pár speciális közülük: \cfn allOf/1, \cfn anyOf/1, \cfn eachOf/1 és az \cfn isnt/1.
Azok, amelyek ,,Of''-al végzőnek egy C4 listát várnak a konkrét paraméterükként, ahol mindegyik elem magukban egy szűrő.
Ezek közül az \cfn allOf/1 esetén mindegyiknek, másik két esetben az egyiknek kell teljesülnie a találat sikerességéhez.
A különbség, hogy az \cfn eachOf/1 esetén a belső értékek közül, mindegyik amelyik sikeres azt kezelő függvény (ha van hozzá csatolva) meg fogja kapni lekezelésre, míg az \cfn anyOf/1 esetén csak az elsőt kezeli.
A nevesített kezelőket \ld az alábbi szakaszban.

Az \cfn isnt/1 függvény negál egy másik szűrőt.
Lehetséges ugyan erre a célra az \cfn (\~)/1 prefix operátort is használni.
Utóbbi listák esetén hasznos lehet, hogy a \cfn (,)/2 operátor precedenciája ne zavarja a megfelelő kiértékelést, úgy mint a függvényeknél.

Végül, a \cfn that/1 függvényt lehetséges lecserélni: alternatívája a \cfn which/1 függvény, ami egy listát vár, amelyiknek elemeiről mind elvárja, hogy igazak legyenek.
Ez különbözik, az \cfn allOf/1 függvénytől annyiban, hogy utóbbi egy saját {\tt "allOf(\dots)"} szűrőt helyez el a lekérdezésben, míg a \cfn which/1 a kapott paramétereit a szülő szűrőnek, mint lista átadja.
Ennek előnye, hogy vannak olyan szűrők, amelyek képesek több értéket értelmezni, mint paraméter.

Az eddigiekben kerültem a kezelő függvények megadását.
Ezek azok a függvények, amelyek a sikeres találatokat a blAST programot futtató számára hasznosíthatóvá teszik a keresést.

Egy utolsó megjegyzést még érdemes megtenni: mivel minden az előzőekben tárgyalt függvény csak normál C4 vagy C4 natív kódként került implementálásra, amennyiben külső kód fenntartja a megfelelő interfészeket, amelyet a többi függvény elvár, szabadon bővíthető a rendszer.
Egy jó példa lehet, egy saját szűrő függvény, mert gyakran szűr valaki arra, hogy egy hívott függvény neve ,,{\tt malloc}'' vagy ,,{\tt calloc}''.
Ehhez alapvetően ki kellene írni, hogy {\tt anyOf (named "malloc"), (named "calloc"), nil}.
De mivel ezek mind csak primitív C4 függvények, így saját szűrőt az alábbi módon lehet könnyen létrehozni, és használni is.
Ez a kód azokat a függvényhívás kifejezéseket keresi meg, amelyek az előbb említett két függvény valamelyikét hívják.
(Az \cfn is/1 függvény opcionális, csak a kifejezések olvashatóságának könnyítését szolgálja, szemantikailag semmit nem csinál.)

\begin{lstlisting}
let memcall \anyOf (named "malloc"), (named "calloc"), nil
# ...
such call expression
that calls
  such function declaration that is memcall
\end{lstlisting}

\subsubsection*{Nevesített találatokat kezelő függvények}

Egyes komplex lekérdezések esetén hasznos a rész elemeket megnevezni.
Ezek előnye, hogy a nagy lekérdezések részeivel is lehetséges találatokat kezelő függvényt rendelni, anélkül, hogy duplikálni kéne annak logikáját egy másik lekérdezésbe.
Ez esetlegesen nem is egyszerűen megvalósítható, mert lehet, hogy csak olyan specifikus esetekben kerül elő, amelyek a nagy lekérdezés kontextusában állnak csak elő.
Az alább látható \aref{lst:complex}.~listában a {``{\tt malloc}''} nevű deklarációk csak akkor kerülnek megtalálásra, ha legalább egy helyen meghívásra kerülnek.
Ezt nem könnyű, ha egyáltalán lehetséges, külön lekérdezésben replikálni.

Az ilyen esetek megoldására lehet elnevezni belső találatokat.
Ezek a ,{\tt declaration/0}' és ,{\tt expression/0}' függvények és két módon lehet nevet rendelni hozzájuk:
\begin{itemize}
\item
A {,{\tt (@\#)/2}'} operátorral, direkt a nevet megadva a találatot bevezető függvény után.
A már említett \aref{lst:complex}.~listában ez látható.

\item
A {,{\tt bound/2}'} függvénnyel, direkt a nevet a találatot bevezető függvény és módosítói előtt.
A már említett \aref{lst:complex}.~listában láthatóval ekvivalens a következő lista:
\begin{lstlisting}
such bound "malloc" function declaration
that is named "malloc"
\end{lstlisting}
\end{itemize}
Ezen két szintaxis teljesen ekvivalens, az operátor a függvénnyel van implementálva.

A blAST keresések \gls{DSL}-jének legfelső szintjén -- ahova a ,{\tt match/2}' függvényeket is kell írni -- meg lehet adni ,,{\tt handle}'' függvényeket.
Ezeken keresztül lehet megadni, hogy az előzőek szerint elnevezett belső találatok esetén, milyen kódrészlet fusson le.
A paraméterezésük a név után azonosa a ,{\tt match/2}' függvény első paraméterével.
Erről \ld az előző szakaszt.

Ha elnevezésre kerül egy belső találat, de nem kerül hozzá kezelő függvény, akkor találatkor egy figyelmeztető diagnosztika megjelenik, hogy nem talált kezelő függvényt az adott nevű találathoz a blAST kiértékelő motorja.
Ennek a formátuma a normális diagnosztikai formátum, amivel a megadott ,,{\tt diagnostic/1}'' kezelők is kiírják a találatokat.

\begin{lstlisting}[caption={Összetett keresések blAST-ban},label={lst:complex}]
blast used_ast
	match into diagnostic "dynalloc_call"
	such call expression
	that calls
		such function declaration@#"malloc"
		that is named "malloc"

	match into diagnostic "global variable defined"
	such variable declaration
	which global
		 , ~static_local
		 , (expanded in main_file)
		 , nil
	
	handle "malloc" with diagnostic "malloc declared here"
end_blast
\end{lstlisting}

